\subsection{محدودیت‌ها و تحلیل شکست}
بررسی ساختار درونی فشرده‌سازی حوزه فرکانس مبین چالش‌های زیر است:
\begin{itemize}
    \item \textbf{پدیده نشت طیفی و لبه‌های ناپیوسته (\lr{Spectral Leakage}):} ماتریس‌های وزن بر خلاف سیگنال‌های تصویری یا صوتی پیوسته دارای دور تناوب طبیعی نیستند. اعمال تبدیل فوریه دو‌بعدی بدون اعمال پنجره‌گذاری (\lr{Windowing}) می‌تواند موجب پراکندگی انرژی در فرکانس‌های بالا شده و خطای تقریب در بازسازی لبه‌های ماتریس وزن را افزایش دهد.
    \item \textbf{کاهش دقت در مدل‌های حساس به فاز (\lr{Phase Distortions}):} فرآیند ماسک‌گذاری ضرایب مختلط و حذف ضرایب ضعیف فاز می‌تواند توزیع فعال‌سازی‌های غیرخطی بعدی را در مدل‌های پیچیده دچار اعوجاج نماید، که افت دقت در مدل \lr{RoBERTa} به \lr{72.02\%} بازتابی از این پدیده است.
    \item \textbf{سربار فاز محاسبه تبدیل‌های فوریه دو‌بعدی:} محاسبه الگوریتم‌های \lr{FFT2} برای ماتریس‌های وزن بسیار بزرگ نیازمند حافظه موقت دوچندان جهت نگهداری تانسورهای مختلط \lr{complex64} در حین فرآیند کالیبراسیون اولیه است.
\end{itemize}
